
Dieses Handbuch beschreibt die Installation, den Aufbau sowie die Bedienung der im Rahmen der Diplomarbeit 
entwickelten VR-Simulation. Die Anwendung wurde mit Unity entwickelt und ermöglicht die 
Perspektive von Rollstuhlfahrenden in einer virtuellen Stadtumgebung einzunehmen.

Ziel der Simulation ist es, typische Hindernisse wie Bankautomaten, Bordsteinkanten und Pflastersteine 
realitätsnah erfahrbar zu machen und deren Auswirkungen auf die Mobilität nachvollziehbar darzustellen.

\section{VR-Welt}
\setauthor{Elisa Kaltenberger}
Zuerst wird Unity installiert und das Projekt geklont. Anschließend wird die entsprechende Szene geöffnet. Falls gewünscht, kann eine VR-Brille angeschlossen und die Szene gestartet werden. Weitere Einstellungen müssen nicht vorgenommen werden.

\subsection{Voraussetzungen}
\setauthor{Elisa Kaltenberger}
Für die Nutzung der Simulation werden folgende Komponenten benötigt:
\begin{itemize}
    \item Laptop oder PC mit installierter Unity-Version (passend zur Projektdatei)
    \item Optional: VR-Headset (z. B. Meta Quest 2 oder vergleichbares Gerät)
    \item Installierte Meta Quest Link-Software bei Verwendung einer Meta-Brille
    \item USB-C-Kabel zur Verbindung zwischen Headset und Computer
\end{itemize}

\subsection{Installation}
\setauthor{Elisa Kaltenberger}
\begin{itemize}
    \item Unity installieren (falls noch nicht vorhanden)
    \item Projektdatei klonen oder herunterladen
    \item Projekt über den Unity Hub öffnen
    \item Die Hauptszene im Projektordner auswählen und öffnen
    \item Optional: VR-Brille anschließen
    \item Auf „Play-Button“ klicken => die Simulation startet
\end{itemize}


\subsection{Ohne VR-Brille}
\setauthor{Elisa Kaltenberger}
Die Simulation kann auch ohne Headset direkt am Bildschirm getestet werden. Die Steuerun erfolgt über 
WSAD-Tasten und Pfeiltasten für die Bewegung und die Maus für die Blickrichtung. Es ist jedoch zu beachten, 
dass die Erfahrung ohne VR-Brille eingeschränkt ist, da die Immersion und das räumliche Bewusstsein reduziert 
sind.

\subsection{Mit VR-Brille}
\setauthor{Elisa Kaltenberger}
Für das vollständige immersive Erlebnis wird die Nutzung mit VR-Headset empfohlen.

\begin{itemize}
    \item Meta Quest Link starten
    \item VR-Brille per Kabel mit dem Laptop verbinden
    \item Warten, bis das Headset vom System erkannt wurde (dies kann einige Minuten dauern)
    \item Szene in Unity starten
\end{itemize}

Steuerung in VR:
\begin{itemize}
    \item Bewegung durch Neigen und Ausrichten der Controller in die gewünschte Richtung
    \item Kopfbewegungen steuern die Blickrichtung
\end{itemize}


